\documentclass[10pt,oneside,openany]{memoir}

\iffalse
%font style 1
\usepackage[p,osf]{scholax}
\usepackage{amsmath,amsthm,amssymb}
\usepackage[scaled=1.075,ncf,vvarbb]{newtxmath}
\fi

%font style 2
\usepackage[T1]{fontenc}
\usepackage{amsmath}
\usepackage[fixamsmath]{mathtools}  % Extension to amsmath
\usepackage{amsthm}
\usepackage{amssymb}
\renewcommand*{\mathbf}[1]{\varmathbb{#1}}
\usepackage{newpxtext}
\usepackage{eulerpx}
\usepackage{mathrsfs}

\usepackage{eucal}
\usepackage{datetime}
    \newdateformat{specialdate}{\THEYEAR\ \monthname\ \THEDAY}
\usepackage[margin=1.5in]{geometry}
\usepackage{fancyhdr}
    \pagestyle{fancy}
    \fancyhf{}
    \cfoot{\scriptsize\thepage}
    \fancyhead[R]{\scalebox{0.8}{\nouppercase{\rightmark}}}
    \fancyhead[L]{\scalebox{0.8}{\nouppercase{\leftmark}}}

    \fancypagestyle{plain}{%
        \fancyhf{}
        \cfoot{\scriptsize\thepage}
        \fancyhead[R]{\scalebox{0.8}{\phantom{a}}}
        \fancyhead[L]{\scalebox{0.8}{Applied Analysis}}
    }
\usepackage{keytheorems}
    \newkeytheoremstyle{def}{
    inherit-style=definition,
    spaceabove=1.5em,
    spacebelow=1.5em,
    postheadspace=0.5em,
    }
    \newkeytheoremstyle{thm}{
    spaceabove=1.5em,
    spacebelow=1.5em,
    postheadspace=0.5em,
    }
    \newkeytheorem{theorem}[
    name = Theorem,
    style=thm,
    numberwithin = section,
    ]
    \newkeytheorem{theorem*}[
    name=Theorem,
    style=thm,
    numbered=false
    ]
    \newkeytheorem{proposition}[
    name = Proposition,
    style=thm,
    numberwithin = section,
    sibling=theorem
    ]
    \newkeytheorem{lemma}[
    name = Lemma,
    style=thm,
    numberwithin = section,
    sibling=theorem
    ]
    \newkeytheorem{corollary}[
    name = Corollary,
    style=thm,
    numberwithin = section,
    sibling=theorem
    ]
    \newkeytheorem{definition}[
    name = Definition,
    style=def,
    numberwithin = section,
    style = definition
    ]
    \newkeytheorem{example}[
    name = Example,
    style=def,
    numberwithin = section,
    style = definition
    ]
    \newkeytheorem{remark}[
    name=Remark,
    style=remark, 
    numbered=false  
    ]
    \newkeytheorem{question}[
    name=Question,
    style=remark, 
    numbered=false  
    ]
    \newkeytheorem{exercise}[
    name = Exercise,
    style=thm
    ]
    \newkeytheorem{axiom}[name = Axiom]
\usepackage{enumitem}
\usepackage[explicit]{titlesec}
    \titleformat{\chapter}[display]
    {\bfseries\LARGE\raggedright}
    {Chapter {\thechapter}}
    {.1ex minus .1ex}
    {\Large #1}
    \titlespacing{\chapter}
    {3pc}{*3}{40pt}[3pc]

    \titleformat{\section}[block]
    {\normalfont\bfseries\large}
    {}{0pt}
    {\fbox{\strut \S\ \thesection.\ #1}}
    \titlespacing{\section}
    {0pt}{3ex plus .1ex minus .2ex}{3ex plus .1ex minus .2ex}
    \setlength{\fboxsep}{4pt}   % padding inside the box
    \setlength{\fboxrule}{0.6pt} % border thickness
    
    \titleformat{name=\section,numberless}[block]
    {\normalfont\bfseries\large\raggedright}
    {}{0pt}
    {\fbox{\strut #1}}
\usepackage[utf8x]{inputenc}
\usepackage{tikz}
\usepackage{tikz-cd}
\usepackage{float}
\usetikzlibrary{patterns}
\usetikzlibrary{positioning,arrows.meta}
\usetikzlibrary{decorations.pathreplacing}
\usepackage{wasysym}
\usepackage{hyperref}
\usepackage{ulem}
\usepackage{pgfplots}
\usetikzlibrary{patterns}
\usepgfplotslibrary{fillbetween}
\pgfplotsset{compat=1.18}

\renewcommand{\headrulewidth}{0pt} % removes the top/header line
% (optional) if you also want no line above the footer:
\renewcommand{\footrulewidth}{0pt}
%%%%%%%%%%%%%%%%%%%%%%%%%%%%%%%%%%%%%%%%%%%%%%%%%%%%%%%%%%%%%
%%%%%%%%%%%%%%%%%%%%%%%%%%%%%%%%%%%%%%%%%%%%%%%%%%%%%%%%%%%%%
% ============================================================
% makros.tex — reorganized, full restored version
% ============================================================

% ============================================================
% Sha / Cyrillic support notes
% ============================================================

%to make the correct symbol for Sha
%\newcommand\cyr{%
%\renewcommand\rmdefault{wncyr}%
%\renewcommand\sfdefault{wncyss}%
%\renewcommand\encodingdefault{OT2}%
%\normalfont \selectfont} \DeclareTextFontCommand{\textcyr}{\cyr}

% ============================================================
% Math operators
% ============================================================

\DeclareMathOperator{\ab}{ab}
\DeclareMathOperator{\ad}{ad}
\DeclareMathOperator{\adj}{adj}
\DeclareMathOperator{\alg}{alg}
\DeclareMathOperator{\Alt}{Alt}
\DeclareMathOperator{\Ann}{Ann}
\DeclareMathOperator{\arith}{arith}
\DeclareMathOperator{\Aut}{Aut}
\DeclareMathOperator{\Be}{B}
\DeclareMathOperator{\Bd}{Bd}
\DeclareMathOperator{\card}{card}
\DeclareMathOperator{\Char}{char}
\DeclareMathOperator{\csp}{csp}
\DeclareMathOperator{\codim}{codim}
\DeclareMathOperator{\coker}{coker}
\DeclareMathOperator{\coh}{H}
\DeclareMathOperator{\compl}{compl}
\DeclareMathOperator{\conj}{conj}
\DeclareMathOperator{\cont}{cont}
\DeclareMathOperator{\Cov}{Cov}
\DeclareMathOperator{\crys}{crys}
\DeclareMathOperator{\Crys}{Crys}
\DeclareMathOperator{\cusp}{cusp}
\DeclareMathOperator{\diag}{diag}
\DeclareMathOperator{\diam}{diam}
\DeclareMathOperator{\Dom}{Dom}
\DeclareMathOperator{\disc}{disc}
\DeclareMathOperator{\dist}{dist}
\DeclareMathOperator{\Div}{div}
\DeclareMathOperator{\dR}{dR}
\DeclareMathOperator{\Eis}{Eis}
\DeclareMathOperator{\End}{End}
\DeclareMathOperator{\ev}{ev}
\DeclareMathOperator{\eval}{eval}
\DeclareMathOperator{\Eq}{Eq}
\DeclareMathOperator{\Ext}{Ext}
\DeclareMathOperator{\Fil}{Fil}
\DeclareMathOperator{\Fitt}{Fitt}
\DeclareMathOperator{\Frob}{Frob}
\DeclareMathOperator{\G}{G}
\DeclareMathOperator{\Gal}{Gal}
\DeclareMathOperator{\GL}{GL}
\DeclareMathOperator{\Gr}{Gr}
\DeclareMathOperator{\Graph}{Graph}
\DeclareMathOperator{\GSp}{GSp}
\DeclareMathOperator{\GUn}{GU}
\DeclareMathOperator{\Hom}{Hom}
\DeclareMathOperator{\id}{id}
\DeclareMathOperator{\Id}{Id}
\DeclareMathOperator{\Ik}{Ik}
\DeclareMathOperator{\IM}{Im}
\DeclareMathOperator{\Image}{im}
\DeclareMathOperator{\Ind}{Ind}
\DeclareMathOperator{\Inf}{inf}
\DeclareMathOperator{\ins}{ins}
\DeclareMathOperator{\inv}{inv}
\DeclareMathOperator{\Isom}{Isom}
\DeclareMathOperator{\J}{J}
\DeclareMathOperator{\Jac}{Jac}
\DeclareMathOperator{\lcm}{lcm}
\DeclareMathOperator{\length}{length}
\DeclareMathOperator*{\limit}{limit}
\DeclareMathOperator{\Log}{Log}
\DeclareMathOperator{\M}{M}
\DeclareMathOperator{\Mat}{Mat}
\DeclareMathOperator{\N}{N}
\DeclareMathOperator{\Nm}{Nm}
\DeclareMathOperator{\NIk}{N-Ik}
\DeclareMathOperator{\NSK}{N-SK}
\DeclareMathOperator{\new}{new}
\DeclareMathOperator{\obj}{obj}
\DeclareMathOperator{\old}{old}
\DeclareMathOperator{\ord}{ord}
\DeclareMathOperator{\Or}{O}
\DeclareMathOperator{\op}{op}
\DeclareMathOperator{\out}{out}
\DeclareMathOperator{\PGL}{PGL}
\DeclareMathOperator{\PGSp}{PGSp}
\DeclareMathOperator{\rank}{rank}
\DeclareMathOperator{\Ran}{Ran}
\DeclareMathOperator{\rel}{Rel}
\DeclareMathOperator{\Real}{Re}
\DeclareMathOperator{\RES}{res}
\DeclareMathOperator{\Res}{Res}
%\DeclareMathOperator{\Sha}{\textcyr{Sh}}
\DeclareMathOperator{\Sel}{Sel}
\DeclareMathOperator{\semi}{ss}
\DeclareMathOperator{\sgn}{sign}
\DeclareMathOperator{\SK}{SK}
\DeclareMathOperator{\SL}{SL}
\DeclareMathOperator{\SO}{SO}
\DeclareMathOperator{\Sp}{Sp}
\DeclareMathOperator{\Span}{span}
\DeclareMathOperator{\Spec}{Spec}
\DeclareMathOperator{\spin}{spin}
\DeclareMathOperator{\st}{st}
\DeclareMathOperator{\St}{St}
\DeclareMathOperator{\SUn}{SU}
\DeclareMathOperator{\supp}{supp}
\DeclareMathOperator{\Sup}{sup}
\DeclareMathOperator{\Sym}{Sym}
\DeclareMathOperator{\Tam}{Tam}
\DeclareMathOperator{\tors}{tors}
\DeclareMathOperator{\tr}{tr}
\DeclareMathOperator{\Tr}{Tr}
\DeclareMathOperator{\un}{un}
\DeclareMathOperator{\Un}{U}
\DeclareMathOperator{\val}{val}
\DeclareMathOperator{\vol}{vol}

% ============================================================
% General aliases and short macros
% ============================================================

\newcommand{\absgal}{\G_{\bbQ}}
\newcommand{\Loc}{L_{\operatorname{loc}}}
\renewcommand{\k}{\kappa}
\newcommand{\Ff}{F_{f}}
%\newcommand{\ts}{\,^{t}\!}
\newcommand{\lat}{\mathscr{L}}
\newcommand{\ov}[1]{\overline{#1}}
\newcommand{\up}{\upsilon}
\newcommand{\e}{\epsilon}
\newcommand{\tac}{\textasteriskcentered}

%mahesh macros
\newcommand{\tm}{\textrm}

\newcommand{\bone}{\mathbf{1}}
\newcommand{\sign}{\mathrm{sign}}
\newcommand{\eps}{\varepsilon}
\newcommand{\textui}[1]{\underline{\textit{#1}}}

%\newcommand{\ov}{\overline}
%\newcommand{\un}{\underline}
\newcommand{\fin}{\mathrm{fin}}
\newcommand{\nl}{\newline \mbox{}}
\newcommand{\h}[1]{\hspace{#1pt}}
\newcommand{\mtext}[1]{\hspace{6pt}\text{#1}\hspace{6pt}}
\newcommand{\lnorm}{\left\lVert}
\newcommand{\rnorm}{\right\rVert}
\newcommand{\ds}{\displaystyle}
\newcommand{\ts}{\textstyle}
\newcommand{\sfrac}[2]{{}^{#1}\mskip -5mu/\mskip -3mu_{#2}}

% ============================================================
% Category-style operators and contoured variants
% ============================================================

\DeclareMathOperator{\Sets}{S \mkern1.04mu e \mkern1.04mu t \mkern1.04mu s}
    \newcommand{\cSets}{\scalebox{1.02}{\contour{black}{$\Sets$}}}
    
\DeclareMathOperator{\Groups}{G \mkern1.04mu r \mkern1.04mu o \mkern1.04mu u \mkern1.04mu p \mkern1.04mu s}
    \newcommand{\cGroups}{\scalebox{1.02}{\contour{black}{$\Groups$}}}

\DeclareMathOperator{\TTop}{T \mkern1.04mu o \mkern1.04mu p}
    \newcommand{\cTop}{\scalebox{1.02}{\contour{black}{$\TTop$}}}

\DeclareMathOperator{\Htp}{H \mkern1.04mu t \mkern1.04mu p}
    \newcommand{\cHtp}{\scalebox{1.02}{\contour{black}{$\Htp$}}}

\DeclareMathOperator{\Mod}{M \mkern1.04mu o \mkern1.04mu d}
    \newcommand{\cMod}{\scalebox{1.02}{\contour{black}{$\Mod$}}}

\DeclareMathOperator{\Ab}{A \mkern1.04mu b}
    \newcommand{\cAb}{\scalebox{1.02}{\contour{black}{$\Ab$}}}

\DeclareMathOperator{\Rings}{R \mkern1.04mu i \mkern1.04mu n \mkern1.04mu g \mkern1.04mu s}
    \newcommand{\cRings}{\scalebox{1.02}{\contour{black}{$\Rings$}}}

\DeclareMathOperator{\ComRings}{C \mkern1.04mu o \mkern1.04mu m \mkern1.04mu R \mkern1.04mu i \mkern1.04mu n \mkern1.04mu g \mkern1.04mu s}
    \newcommand{\cComRings}{\scalebox{1.05}{\contour{black}{$\ComRings$}}}

\DeclareMathOperator{\hHom}{H \mkern1.04mu o \mkern1.04mu m}
    \newcommand{\cHom}{\scalebox{1.02}{\contour{black}{$\hHom$}}}

% ============================================================
% Mathcal
% ============================================================

\newcommand{\cA}{\mathcal{A}}
\newcommand{\cB}{\mathcal{B}}
\newcommand{\cC}{\mathcal{C}}
\newcommand{\cD}{\mathcal{D}}
\newcommand{\cE}{\mathcal{E}}
\newcommand{\cF}{\mathcal{F}}
\newcommand{\cG}{\mathcal{G}}
\newcommand{\cH}{\mathcal{H}}
\newcommand{\cI}{\mathcal{I}}
\newcommand{\cJ}{\mathcal{J}}
\newcommand{\cK}{\mathcal{K}}
\newcommand{\cL}{\mathcal{L}}
\newcommand{\cM}{\mathcal{M}}
\newcommand{\cN}{\mathcal{N}}
\newcommand{\cO}{\mathcal{O}}
\newcommand{\cP}{\mathcal{P}}
\newcommand{\cQ}{\mathcal{Q}}
\newcommand{\cR}{\mathcal{R}}
\newcommand{\cS}{\mathcal{S}}
\newcommand{\cT}{\mathcal{T}}
\newcommand{\cU}{\mathcal{U}}
\newcommand{\cV}{\mathcal{V}}
\newcommand{\cW}{\mathcal{W}}
\newcommand{\cX}{\mathcal{X}}
\newcommand{\cY}{\mathcal{Y}}
\newcommand{\cZ}{\mathcal{Z}}

% ============================================================
% Mathscrl% 
%============================================================

\newcommand{\scra}{\mathscr{a}}
\newcommand{\scrA}{\mathscr{A}}
\newcommand{\scrb}{\mathscr{b}}
\newcommand{\scrB}{\mathscr{B}}
\newcommand{\scrc}{\mathscr{c}}
\newcommand{\scrC}{\mathscr{C}}
\newcommand{\scrd}{\mathscr{d}}
\newcommand{\scrD}{\mathscr{D}}
\newcommand{\scre}{\mathscr{e}}
\newcommand{\scrE}{\mathscr{E}}
\newcommand{\scrf}{\mathscr{f}}
\newcommand{\scrF}{\mathscr{F}}
\newcommand{\scrg}{\mathscr{g}}
\newcommand{\scrG}{\mathscr{G}}
\newcommand{\scrh}{\mathscr{h}}
\newcommand{\scrH}{\mathscr{H}}
\newcommand{\scri}{\mathscr{i}}
\newcommand{\scrI}{\mathscr{I}}
\newcommand{\scrj}{\mathscr{j}}
\newcommand{\scrJ}{\mathscr{J}}
\newcommand{\scrk}{\mathscr{k}}
\newcommand{\scrK}{\mathscr{K}}
\newcommand{\scrl}{\mathscr{l}}
\newcommand{\scrL}{\mathscr{L}}
\newcommand{\scrm}{\mathscr{m}}
\newcommand{\scrM}{\mathscr{M}}
\newcommand{\scrn}{\mathscr{n}}
\newcommand{\scrN}{\mathscr{N}}
\newcommand{\scro}{\mathscr{o}}
\newcommand{\scrO}{\mathscr{O}}
\newcommand{\scrp}{\mathscr{p}}
\newcommand{\scrP}{\mathscr{P}}
\newcommand{\scrq}{\mathscr{q}}
\newcommand{\scrQ}{\mathscr{Q}}
\newcommand{\scrr}{\mathscr{r}}
\newcommand{\scrR}{\mathscr{R}}
\newcommand{\scrs}{\mathscr{s}}
\newcommand{\scrS}{\mathscr{S}}
\newcommand{\scrt}{\mathscr{t}}
\newcommand{\scrT}{\mathscr{T}}
\newcommand{\scru}{\mathscr{u}}
\newcommand{\scrU}{\mathscr{U}}
\newcommand{\scrv}{\mathscr{v}}
\newcommand{\scrV}{\mathscr{V}}
\newcommand{\scrw}{\mathscr{w}}
\newcommand{\scrW}{\mathscr{W}}
\newcommand{\scrx}{\mathscr{x}}
\newcommand{\scrX}{\mathscr{X}}
\newcommand{\scry}{\mathscr{y}}
\newcommand{\scrY}{\mathscr{Y}}
\newcommand{\scrz}{\mathscr{z}}
\newcommand{\scrZ}{\mathscr{Z}}

% ============================================================
% Mathfrak (missing \fi)
% ============================================================

\newcommand{\fa}{\mathfrak{a}}
\newcommand{\fA}{\mathfrak{A}}
\newcommand{\fb}{\mathfrak{b}}
\newcommand{\fB}{\mathfrak{B}}
\newcommand{\fc}{\mathfrak{c}}
\newcommand{\fC}{\mathfrak{C}}
\newcommand{\fd}{\mathfrak{d}}
\newcommand{\fD}{\mathfrak{D}}
\newcommand{\fe}{\mathfrak{e}}
\newcommand{\fE}{\mathfrak{E}}
\newcommand{\ff}{\mathfrak{f}}
\newcommand{\fF}{\mathfrak{F}}
\newcommand{\fg}{\mathfrak{g}}
\newcommand{\fG}{\mathfrak{G}}
\newcommand{\fh}{\mathfrak{h}}
\newcommand{\fH}{\mathfrak{H}}
\newcommand{\fI}{\mathfrak{I}}
\newcommand{\fj}{\mathfrak{j}}
\newcommand{\fJ}{\mathfrak{J}}
\newcommand{\fk}{\mathfrak{k}}
\newcommand{\fK}{\mathfrak{K}}
\newcommand{\fl}{\mathfrak{l}}
\newcommand{\fL}{\mathfrak{L}}
\newcommand{\fm}{\mathfrak{m}}
\newcommand{\fM}{\mathfrak{M}}
\newcommand{\fn}{\mathfrak{n}}
\newcommand{\fN}{\mathfrak{N}}
\newcommand{\fo}{\mathfrak{o}}
\newcommand{\fO}{\mathfrak{O}}
\newcommand{\fp}{\mathfrak{p}}
\newcommand{\fP}{\mathfrak{P}}
\newcommand{\fq}{\mathfrak{q}}
\newcommand{\fQ}{\mathfrak{Q}}
\newcommand{\fr}{\mathfrak{r}}
\newcommand{\fR}{\mathfrak{R}}
\newcommand{\fs}{\mathfrak{s}}
\newcommand{\fS}{\mathfrak{S}}
\newcommand{\ft}{\mathfrak{t}}
\newcommand{\fT}{\mathfrak{T}}
\newcommand{\fu}{\mathfrak{u}}
\newcommand{\fU}{\mathfrak{U}}
\newcommand{\fv}{\mathfrak{v}}
\newcommand{\fV}{\mathfrak{V}}
\newcommand{\fw}{\mathfrak{w}}
\newcommand{\fW}{\mathfrak{W}}
\newcommand{\fx}{\mathfrak{x}}
\newcommand{\fX}{\mathfrak{X}}
\newcommand{\fy}{\mathfrak{y}}
\newcommand{\fY}{\mathfrak{Y}}
\newcommand{\fz}{\mathfrak{z}}
\newcommand{\fZ}{\mathfrak{Z}}

% ============================================================
% Mathbf
% ============================================================

\newcommand{\bfA}{\mathbf{A}}
\newcommand{\bfB}{\mathbf{B}}
\newcommand{\bfC}{\mathbf{C}}
\newcommand{\bfD}{\mathbf{D}}
\newcommand{\bfE}{\mathbf{E}}
\newcommand{\bfF}{\mathbf{F}}
\newcommand{\bfG}{\mathbf{G}}
\newcommand{\bfH}{\mathbf{H}}
\newcommand{\bfI}{\mathbf{I}}
\newcommand{\bfJ}{\mathbf{J}}
\newcommand{\bfK}{\mathbf{K}}
\newcommand{\bfL}{\mathbf{L}}
\newcommand{\bfM}{\mathbf{M}}
\newcommand{\bfN}{\mathbf{N}}
\newcommand{\bfO}{\mathbf{O}}
\newcommand{\bfP}{\mathbf{P}}
\newcommand{\bfQ}{\mathbf{Q}}
\newcommand{\bfR}{\mathbf{R}}
\newcommand{\bfS}{\mathbf{S}}
\newcommand{\bfT}{\mathbf{T}}
\newcommand{\bfU}{\mathbf{U}}
\newcommand{\bfV}{\mathbf{V}}
\newcommand{\bfW}{\mathbf{W}}
\newcommand{\bfX}{\mathbf{X}}
\newcommand{\bfY}{\mathbf{Y}}
\newcommand{\bfZ}{\mathbf{Z}}

\newcommand{\bfa}{\mathbf{a}}
\newcommand{\bfb}{\mathbf{b}}
\newcommand{\bfc}{\mathbf{c}}
\newcommand{\bfd}{\mathbf{d}}
\newcommand{\bfe}{\mathbf{e}}
\newcommand{\bff}{\mathbf{f}}
\newcommand{\bfg}{\mathbf{g}}
\newcommand{\bfh}{\mathbf{h}}
\newcommand{\bfi}{\mathbf{i}}
\newcommand{\bfj}{\mathbf{j}}
\newcommand{\bfk}{\mathbf{k}}
\newcommand{\bfl}{\mathbf{l}}
\newcommand{\bfm}{\mathbf{m}}
\newcommand{\bfn}{\mathbf{n}}
\newcommand{\bfo}{\mathbf{o}}
\newcommand{\bfp}{\mathbf{p}}
\newcommand{\bfq}{\mathbf{q}}
\newcommand{\bfr}{\mathbf{r}}
\newcommand{\bfs}{\mathbf{s}}
\newcommand{\bft}{\mathbf{t}}
\newcommand{\bfu}{\mathbf{u}}
\newcommand{\bfv}{\mathbf{v}}
\newcommand{\bfw}{\mathbf{w}}
\newcommand{\bfx}{\mathbf{x}}
\newcommand{\bfy}{\mathbf{y}}
\newcommand{\bfz}{\mathbf{z}}

% ============================================================
% Blackboard bold
% ============================================================

\newcommand{\bbA}{\mathbb{A}}
\newcommand{\bbB}{\mathbb{B}}
\newcommand{\bbC}{\mathbb{C}}
\newcommand{\bbD}{\mathbb{D}}
\newcommand{\bbE}{\mathbb{E}}
\newcommand{\bbF}{\mathbb{F}}
\newcommand{\bbG}{\mathbb{G}}
\newcommand{\bbH}{\mathbb{H}}
\newcommand{\bbI}{\mathbb{I}}
\newcommand{\bbJ}{\mathbb{J}}
\newcommand{\bbK}{\mathbb{K}}
\newcommand{\bbL}{\mathbb{L}}
\newcommand{\bbM}{\mathbb{M}}
\newcommand{\bbN}{\mathbb{N}}
\newcommand{\bbO}{\mathbb{O}}
\newcommand{\bbP}{\mathbb{P}}
\newcommand{\bbQ}{\mathbb{Q}}
\newcommand{\bbR}{\mathbb{R}}
\newcommand{\bbS}{\mathbb{S}}
\newcommand{\bbT}{\mathbb{T}}
\newcommand{\bbU}{\mathbb{U}}
\newcommand{\bbV}{\mathbb{V}}
\newcommand{\bbW}{\mathbb{W}}
\newcommand{\bbX}{\mathbb{X}}
\newcommand{\bbY}{\mathbb{Y}}
\newcommand{\bbZ}{\mathbb{Z}}

\newcommand{\Bmu}{\mbox{$\raisebox{-0.59ex}
  {$l$}\hspace{-0.18em}\mu\hspace{-0.88em}\raisebox{-0.98ex}{\scalebox{2}
  {$\color{white}.$}}\hspace{-0.416em}\raisebox{+0.88ex}
  {$\color{white}.$}\hspace{0.46em}$}{}}  %need graphicx and xcolor. this produces blackboard bold mu 

% ============================================================
% Greek-letter aliases
% ============================================================

\newcommand{\jota}{\jmath}

% ============================================================
% Matrices
% ============================================================

\newcommand{\bmat}{\left( \begin{matrix}}
\newcommand{\emat}{\end{matrix} \right)}

\newcommand{\bbmat}{\left[ \begin{matrix}}
\newcommand{\ebmat}{\end{matrix} \right]}

\newcommand{\pmat}{\left( \begin{smallmatrix}}
\newcommand{\epmat}{\end{smallmatrix} \right)}

\newcommand{\mat}[4]{\begin{pmatrix}{#1}&{#2}\\{#3}&{#4}\end{pmatrix}}

% ============================================================
% Residues and averaged integrals
% ============================================================

\newcommand{\res}[1]{\underset{#1}{\RES}\,}

\makeatletter
\newcommand*{\mint}[1]{%
  % #1: overlay symbol
  \mint@l{#1}{}%
}
\newcommand*{\mint@l}[2]{%
  % #1: overlay symbol
  % #2: limits
  \@ifnextchar\limits{%
    \mint@l{#1}%
  }{%
    \@ifnextchar\nolimits{%
      \mint@l{#1}%
    }{%
      \@ifnextchar\displaylimits{%
        \mint@l{#1}%
      }{%
        \mint@s{#2}{#1}%
      }%
    }%
  }%
}
\newcommand*{\mint@s}[2]{%
  % #1: limits
  % #2: overlay symbol
  \@ifnextchar_{%
    \mint@sub{#1}{#2}%
  }{%
    \@ifnextchar^{%
      \mint@sup{#1}{#2}%
    }{%
      \mint@{#1}{#2}{}{}%
    }%
  }%
}
\def\mint@sub#1#2_#3{%
  \@ifnextchar^{%
    \mint@sub@sup{#1}{#2}{#3}%
  }{%
    \mint@{#1}{#2}{#3}{}%
  }%
}
\def\mint@sup#1#2^#3{%
  \@ifnextchar_{%
    \mint@sup@sub{#1}{#2}{#3}%
  }{%
    \mint@{#1}{#2}{}{#3}%
  }%
}
\def\mint@sub@sup#1#2#3^#4{%
  \mint@{#1}{#2}{#3}{#4}%
}
\def\mint@sup@sub#1#2#3_#4{%
  \mint@{#1}{#2}{#4}{#3}%
}
\newcommand*{\mint@}[4]{%
  % #1: \limits, \nolimits, \displaylimits
  % #2: overlay symbol: -, =, ...
  % #3: subscript
  % #4: superscript
  \mathop{}%
  \mkern-\thinmuskip
  \mathchoice{%
    \mint@@{#1}{#2}{#3}{#4}%
        \displaystyle\textstyle\scriptstyle
  }{%
    \mint@@{#1}{#2}{#3}{#4}%
        \textstyle\scriptstyle\scriptstyle
  }{%
    \mint@@{#1}{#2}{#3}{#4}%
        \scriptstyle\scriptscriptstyle\scriptscriptstyle
  }{%
    \mint@@{#1}{#2}{#3}{#4}%
        \scriptscriptstyle\scriptscriptstyle\scriptscriptstyle
  }%
  \mkern-\thinmuskip
  \int#1%
  \ifx\\#3\\\else_{#3}\fi
  \ifx\\#4\\\else^{#4}\fi  
}
\newcommand*{\mint@@}[7]{%
  % #1: limits
  % #2: overlay symbol
  % #3: subscript
  % #4: superscript
  % #5: math style
  % #6: math style for overlay symbol
  % #7: math style for subscript/superscript
  \begingroup
    \sbox0{$#5\int\m@th$}%
    \sbox2{$#5\int_{}\m@th$}%
    \dimen2=\wd0 %
    % => \dimen2 = width of \int
    \let\mint@limits=#1\relax
    \ifx\mint@limits\relax
      \sbox4{$#5\int_{\kern1sp}^{\kern1sp}\m@th$}%
      \ifdim\wd4>\wd2 %
        \let\mint@limits=\nolimits
      \else
        \let\mint@limits=\limits
      \fi
    \fi
    \ifx\mint@limits\displaylimits
      \ifx#5\displaystyle
        \let\mint@limits=\limits
      \fi
    \fi
    \ifx\mint@limits\limits
      \sbox0{$#7#3\m@th$}%
      \sbox2{$#7#4\m@th$}%
      \ifdim\wd0>\dimen2 %
        \dimen2=\wd0 %
      \fi
      \ifdim\wd2>\dimen2 %
        \dimen2=\wd2 %
      \fi
    \fi
    \rlap{%
      $#5%
        \vcenter{%
          \hbox to\dimen2{%
            \hss
            $#6{#2}\m@th$%
            \hss
          }%
        }%
      $%
    }%
  \endgroup
}
\newcommand{\avg}{\mint{-}}
\makeatother

% ============================================================
% Comments and drafting helpers
% ============================================================

%Comments
\newcommand{\com}[1]{\vspace{5 mm}\par \noindent
\marginpar{\textsc{Comment}} \framebox{\begin{minipage}[c]{0.95
\textwidth} \tt #1 \end{minipage}}\vspace{5 mm}\par}

% ============================================================
% Arrows
% ============================================================

\newcommand{\hooktwoheadrightarrow}{%
  \hookrightarrow\mathrel{\mspace{-15mu}}\rightarrow
}

\makeatletter
\newcommand{\xhooktwoheadrightarrow}[2][]{%
  \lhook\joinrel
  \ext@arrow 0359\rightarrowfill@ {#1}{#2}%
  \mathrel{\mspace{-15mu}}\rightarrow
}
\makeatother

\makeatletter
\newcommand*{\da@rightarrow}{\mathchar"0\hexnumber@\symAMSa 4B }
\newcommand*{\da@leftarrow}{\mathchar"0\hexnumber@\symAMSa 4C }
\newcommand*{\xdashrightarrow}[2][]{%
  \mathrel{%
    \mathpalette{\da@xarrow{#1}{#2}{}\da@rightarrow{\,}{}}{}%
  }%
}
\newcommand{\xdashleftarrow}[2][]{%
  \mathrel{%
    \mathpalette{\da@xarrow{#1}{#2}\da@leftarrow{}{}{\,}}{}%
  }%
}
\newcommand*{\da@xarrow}[7]{%
  % #1: below
  % #2: above
  % #3: arrow left
  % #4: arrow right
  % #5: space left 
  % #6: space right
  % #7: math style 
  \sbox0{$\ifx#7\scriptstyle\scriptscriptstyle\else\scriptstyle\fi#5#1#6\m@th$}%
  \sbox2{$\ifx#7\scriptstyle\scriptscriptstyle\else\scriptstyle\fi#5#2#6\m@th$}%
  \sbox4{$#7\dabar@\m@th$}%
  \dimen@=\wd0 %
  \ifdim\wd2 >\dimen@
    \dimen@=\wd2 %   
  \fi
  \count@=2 %
  \def\da@bars{\dabar@\dabar@}%
  \@whiledim\count@\wd4<\dimen@\do{%
    \advance\count@\@ne
    \expandafter\def\expandafter\da@bars\expandafter{%
      \da@bars
      \dabar@ 
    }%
  }%  
  \mathrel{#3}%
  \mathrel{%   
    \mathop{\da@bars}\limits
    \ifx\\#1\\%
    \else
      _{\copy0}%
    \fi
    \ifx\\#2\\%
    \else
      ^{\copy2}%
    \fi
  }%   
  \mathrel{#4}%
}
\makeatother

% ============================================================
% Relation and delimiter spacing
% ============================================================

\renewcommand{\geq}{\geqslant}
\renewcommand{\leq}{\leqslant}
\newcommand{\midd}{\hspace{4pt}\middle|\hspace{4pt}}

% ============================================================
% Left Right Updated Deliminators
% ============================================================

\makeatletter
\NewDocumentCommand\Left{O{}}{%
    \IfValueTF{#1}{\notblank{#1}{\mathopen\bBigg@{#1}}{\left}}{\left}}
\NewDocumentCommand\Right{O{}}{%
    \IfValueTF{#1}{\notblank{#1}{\mathopen\bBigg@{#1}}{\right}}{\right}}
\makeatother

% ============================================================
% Restrictions
% ============================================================

\newcommand{\littletaller}{\mathchoice{\vphantom{\big|}}{}{}{}}
\newcommand\restr[2]{{% we make the whole thing an ordinary symbol
  \left.\kern-\nulldelimiterspace % automatically resize the bar with \right
  #1 % the function
  \littletaller % pretend it's a little taller at normal size
  \right|_{#2} % this is the delimiter
  }}

% ============================================================
% Chapter title formatting
% ============================================================

\newcommand{\chnum}{\titleformat
{\chapter} % command
[display] % shape
{\centering} % format
{\Huge \color{black} \shadowbox{\thechapter}} % label
{-0.5em} % sep (space between the number and title)
{\LARGE \color{black} \underline} % before-code
}

\newcommand{\chunnum}{\titleformat
{\chapter} % command
[display] % shape
{} % format
{} % label
{0em} % sep
{ \begin{flushright} \begin{tabular}{r}  \Huge \color{black}
} % before-code
[
\end{tabular} \end{flushright} \normalsize
] % after-code
}



\begin{document}
\begin{center}
    {\Large Math 540 \\[0.1in]Homework \#4}\\[.175in]
    {Name:} {\underline{Gianluca Crescenzo\hspace*{2in}}}\\[0.15in]
    \end{center}
    \vspace{4pt}
%%%%%%%%%%%%%%%%%%%%%%%%%%%%%%%%%%%%%%%%%%%%%%%%%%
%%%%%%%%%%%%%%%%%%%%%%%%%%%%%%%%%%%%%%%%%%%%%%%%%%
%%%%%%%%%%%%%%%%%%%%%%%%%%%%%%%%%%%%%%%%%%%%%%%%%%
%%%%%%%%%%%%%%%%%%%%%%%%%%%%%%%%%%%%%%%%%%%%%%%%%%
\section*{\S 23 Homework Problems}
\begin{exercise}[manual-num={4}]
Show that if $X$ is an infinite set, it is connected in the finite complement topology.
\end{exercise}
    \begin{proof}
        Suppose towards contradiction there exists a separation of $X$; i.e., we can find $U,V \subseteq X$ open such that $U,V \neq \emptyset$, $U \cap V = \emptyset$, and $U \cup V = X$. Since $U,V$ are open in $X$ and non-empty, it must be the case that their complements are countable. However, note that we can write $X = (U \cap V)^c = U^c \cup V^c$. This is a contradiction, as an infinite set cannot be written as the union of two countable sets.
    \end{proof}
%%%%%%%%%%%%%%%%%%%%%%%%%%%%%%%%%%%%%%%%%%%%%%%%%%
\begin{exercise}[manual-num={5}]
Show that if $X$ has the discrete topology, then $X$ is totally disconnected. Does the converse hold?
\end{exercise}
    \begin{proof}
        Note that for any $a,b \in X$, the set $\{a,b\}$ is not connected as $\{a\},\{b\}$ form a separation. Let $x \in X$ be arbitrary. Suppose there exists a separation $U,V \in \cP(X)$ such that $U,V \neq \emptyset$, $U \cup V = \{x\}$, $\overline{U} \cap V = \emptyset$, and $U \cap \overline{V} = \emptyset$. It must be the case that $U = V = \{x\}$. But $\overline{\{x\}} \cap \{x\} = \{x\} \cap \{x\} = \{x\}$, which is a contradiction. Thus the singletons are the only connected subsets of $X$.

        The converse does not hold. Consider $\bfQ$ as a subspace of $\bfR$. The above argument holds identically for elements in $\bfQ$, however the subspace topology on $\bfQ$ inherited from $\bfR$ is not the discrete topology.
    \end{proof}

\begin{exercise}[manual-num={9}]
Let $A$ be a proper subset of $X$, and let $B$ be a proper subset of $Y$. If $X$ and $Y$ are connected, show that $(X \times Y) \setminus (A \times B)$ is connected.
\end{exercise}
    \begin{proof}
        Note that for every $x \in X \setminus A$, the set $S_{x} := \{x\} \times Y$ is connected, and for every $y \in Y \setminus B$, the set $T_{y} := X \times \{y\}$ is connected. For a fixed $x \in X \setminus A$, note that $S_x$ intersects $T_y$. Hence $U_x:= S_x \cup \left( \bigcup_{y \in Y \setminus B} T_y\right)$ is also connected. Now for two different $x_1,x_2 \in X$, the sets $U_{x_1}$ and $U_{x_2}$ intersect. This means $\bigcup_{x \in X \setminus A}U_x$ is connected. Thus $(X \times Y) \setminus (A \times B)$ is connected, as $\bigcup_{x \in X \setminus A}U_x = (X \times Y) \setminus (A \times B)$.
    \end{proof}
%%%%%%%%%%%%%%%%%%%%%%%%%%%%%%%%%%%%%%%%%%%%%%%%%%
%%%%%%%%%%%%%%%%%%%%%%%%%%%%%%%%%%%%%%%%%%%%%%%%%%
%%%%%%%%%%%%%%%%%%%%%%%%%%%%%%%%%%%%%%%%%%%%%%%%%%
%%%%%%%%%%%%%%%%%%%%%%%%%%%%%%%%%%%%%%%%%%%%%%%%%%
\newpage
\section*{\S 24 Homework Problems}
\begin{exercise}[manual-num={1(c)}]
Show that $\bfR^n$ and $\bfR$ are not homeomorphic if $n > 1$.
\end{exercise}
    \begin{proof}
        Suppose there exists a homeomorphism $f:\bfR \rightarrow \bfR^n$. Let $x \in \bfR$. Then $f':\bfR\setminus\{x\} \rightarrow \bfR^n \setminus \{f(x)\}$ is also a homeomorphism. But $\bfR^n \setminus\{f(x)\}$ is path-connected, whilst $\bfR \setminus\{x\}$ is not, which is a contradiction.
    \end{proof}

\begin{exercise}[manual-num={3}]
Let $f:X \rightarrow X$ be continuous. Show that if $X = [0,1]$, there is a point $x$ such that $f(x) = x$. The point $x$ is called a \textbf{fixed point} of $f$. What happens if $X$ equals $[0,1)$ or $(0,1)$?
\end{exercise}
    \begin{proof}
        Suppose $0$ and $1$ are not fixed points. Define $g:[0,1] \rightarrow [-1,1]$ by $g(x) = f(x) - x$. Then $g(0) = f(0)$ and $g(1) = f(1) - 1$. Since $0$ and $1$ are not fixed points, it must be the case that $g(0) = f(0) > 0$ and $g(1) = f(1) - 1 < 0$. By the Intermediate Value Theorem, there exists an $r \in [0,1]$ such that $g(r) = 0$. Hence $f(r) - r = 0$; i.e., $f(r) = r$.

        Define $f_1:[0,1) \rightarrow [0,1)$ by $f_1(x) = \frac{x+1}{2}$. This map is continuous, but does not have a fixed point. Define $f_2:(0,1) \rightarrow (0,1)$ by $f_2(x) = x^2$. This map is continuous, but does not have a fixed point.
    \end{proof}

\begin{exercise}[manual-num={8}]
    \phantom{a}
    \begin{enumerate}[label = (\alph*),itemsep=1pt,topsep=3pt]
        \item Is a product of path-connected spaces necessarily path-connected?
        \item If $A \subset X$ and $A$ is path-connected, is $\overline{A}$ necessarily path-connected?
    \end{enumerate}
\end{exercise}
    \begin{proof}
        (a) Let $X,Y$ be path-connected topological spaces. Let $(a,b),(c,d) \in X \times Y$. Since $X$ is path-connected, there exists a continuous function $f:[0,1] \rightarrow X$ such that $f(0) = a$ and $f(1) = c$. Since $Y$ is path-connected, there exists a continuous function $g:[0,1] \rightarrow Y$ such that $g(0) = b$ and $g(1) = d$. Define $h:[0,1] \rightarrow X \times Y$ by $h(x) = (f(x),g(x))$. This map is continuous (Theorem 18.4), and $h(0) = (a,b)$ and $h(1) = (c,d)$. Thus $X \times Y$ is path-connected.

        (b) No. Consider the topologist's sine curve $S := \{(x, \sin(\frac{1}{x})) \mid x \in (0,1]\}$. This set is clearly path connected, but $\overline{S} = S \cup \{(0,y) \mid y \in [-1,1]\}$ is not path-connected.
    \end{proof}
%%%%%%%%%%%%%%%%%%%%%%%%%%%%%%%%%%%%%%%%%%%%%%%%%%
%%%%%%%%%%%%%%%%%%%%%%%%%%%%%%%%%%%%%%%%%%%%%%%%%%
%%%%%%%%%%%%%%%%%%%%%%%%%%%%%%%%%%%%%%%%%%%%%%%%%%
%%%%%%%%%%%%%%%%%%%%%%%%%%%%%%%%%%%%%%%%%%%%%%%%%%
\section*{\S 25 Homework Problems}

\begin{exercise}[manual-num={4}]
Let $X$ be locally path connected. Show that every connected open set in $X$ is path connected.
\end{exercise}
    \begin{proof}
        Let $U \subseteq X$ be an open and connected set. Fix $x_0 \in U$ and define:
            \begin{equation*}
            \begin{split}
                A := \{y \in U \mid \text{ there is a path from $y$ to $x_0$ }\}.
            \end{split}
            \end{equation*}
        Give $A$ the subspace topology inherited from $U$. Note that $A \neq \emptyset$ since $x_0 \in A$.

        Let $y \in A$. Then $y \in U$ by inclusion. Since $X$ is locally path connected and since $U$ is open, there exists an open and path connected set $V$ such that $y \in V \subseteq U$. Let $z \in V$. Then there exists a path from $z$ to $y$. But since there exists a path from $y$ to $x_0$, there must exist a path from $z$ to $x_0$; i.e., $z \in A$. Hence $y \in V \subseteq A$, implying that $A$ is open.

        Now let $y \in U \setminus A$. Then $y \in U$ by inclusion. Since $X$ is locally path connected and since $U$ is open, there exists an open and path connected set $V$ such that $y \in V \subseteq U$. Suppose towards contradiction we can find some $z \in V \cap A$. Since $z \in V$, there is a path between $z$ and $y$. Since $z \in A$, there is a path from $z$ to $x_0$. But this means there is a path from $x_0$ to $y$, contradicting the fact that $y \in U \setminus A$. It must be the case that $V \cap A = \emptyset$; i.e., $V \subseteq U \setminus A$. Hence $y \in V \subseteq U \setminus A$, establishing $U\setminus A$ as an open set.

        Since $A$ and $U\setminus A$ are open in $U$, and since $U$ is connected, it must be the case that $A = U$. Thus $U$ is path connected.
    \end{proof}

\begin{exercise}[manual-num={5}]
Let $X$ denote the rational points of the interval $[0,1] \times \{0\}$ of $\bfR^2$. Let $T$ denote the union of all line segments joining the point $P = \{0\} \times \{1\}$ to points of $X$.
\begin{enumerate}[label = (\alph*),itemsep=1pt,topsep=3pt]
    \item Show that $T$ is path connected, but is locally connected only at the point $p$.
    \item Find a subset of $\bfR^2$ that is path connected but is locally connected at none of its points.
\end{enumerate}
\end{exercise}
    \begin{proof}
        (a) Let $x,y \in T$. If $x$ and $y$ lie on the same "ray" towards a rational point on $(\bfQ \cap [0,1]) \times \{0\}$, then surely $x$ and $y$ are path connected. If $x$ and $y$ lie on two different rays, then since $x$ is path connected to $p$ and $p$ is path connected to $y$, it must be that $x$ is path connected to $y$. Hence $T$ is path connected. Clearly $T$ is locally connected at $p$ (this can be seen visually). Let $x \in T$ be arbitrary. Let $\epsilon = \frac{d(x,p)}{2}$. Then not only does $B(x,\epsilon)$ contain the ray which $x$ lies on, but by the density of $\bfQ$ it contains another ray which $x$ does not lie on. As subsets of $B(x,\epsilon)$, these rays are clearly not connected, hence $T$ is locally connected only at the point $p$.
            \begin{center}
            \scalebox{0.8}{
            \begin{tikzpicture}[x=7cm,y=7cm]
            % Apex point P=(0,1)
            \coordinate (P) at (0,1);
            \fill (P) circle (0.8pt) node[above left] {$P$};

            \fill (0,0) circle (0.6pt) node[below left] {$(0,0)$};
            \fill (1,0) circle (0.6pt) node[below] {$(1,0)$};

            \draw[black!60] (P) -- (0,0);
            \draw[black!60] (P) -- (1/10,0);
            \draw[black!60] (P) -- (3/18,0);
            \draw[black!60] (P) -- (1,0);
            \draw[black!60] (P) -- (1/2,0);
            \draw[black!60] (P) -- (1/3,0);
            \draw[black!60] (P) -- (2/3,0);
            \draw[black!60] (P) -- (1/4,0);
            \draw[black!60] (P) -- (3/4,0);
            \draw[black!60] (P) -- (2/5,0);
            \draw[black!60] (P) -- (3/5,0);
            \draw[black!60] (P) -- (10/11,0);
            \draw[black!60] (P) -- (4.2/5,0);

            \node[black!60!black] at (.6,.6) {$T$};
            \end{tikzpicture}}
            \end{center}

        (b) Consider the set $\bfZ \times (\bfQ \cap [0,1])$ as a subset of $\bfR^2$. Let $T_z$ denote the union of all line segments joining the point $\{z\} \times \{0\}$ to the points in $\{z+1\} \times (\bfQ \cap [0,1])$. Consider the set $T := \bigcup_{z \in \bfZ}T_z$.

        
        \begin{center}
        \vspace{10pt}
        \scalebox{0.9}{
        \begin{tikzpicture}[x=4.2cm,y=4.2cm]
        % Base points {z}x{0}
        \coordinate (Zneg) at (-1,0);
        \coordinate (Z0)   at (0,0);
        \coordinate (Z1)   at (1,0);

        \fill (Zneg) circle (0.8pt) node[below] {$(-1,0)$};
        \fill (Z0)   circle (0.8pt) node[below] {$(0,0)$};
        \fill (Z1)   circle (0.8pt);

        % x=0
        \fill (0,1)      circle (0.7pt) node[above] {$(0,1)$};
        % x=1
        \fill (1,0)      circle (0.7pt) node[below] {$(1,0)$};
        \fill (1,1)      circle (0.7pt) node[above] {$(1,1)$};
        % x=2
        \fill (2,0)      circle (0.7pt) node[below] {$(2,0)$};
        \fill (2,1)      circle (0.7pt) node[above] {$(2,1)$};

        % T_{-1}: from (-1,0) to {0}x(Q cap [0,1]) (sampled)
        \draw[black!60] (Zneg) -- (0,0); 
        \draw[black!60] (Zneg) -- (0,1/10);
        \draw[black!60] (Zneg) -- (0,1/4);
        \draw[black!60] (Zneg) -- (0,1/3);
        \draw[black!60] (Zneg) -- (0,1/2);
        \draw[black!60] (Zneg) -- (0,2/3);
        \draw[black!60] (Zneg) -- (0,3/4);
        \draw[black!60] (Zneg) -- (0,4/5);
        \draw[black!60] (Zneg) -- (0,10/11);
        \draw[black!60] (Zneg) -- (0,1);

        % T_{0}: from (0,0) to {1}x(Q cap [0,1]) (sampled)
        \draw[black!60] (Z0) -- (1,0);
        \draw[black!60] (Z0) -- (1,1/10);
        \draw[black!60] (Z0) -- (1,1/4);
        \draw[black!60] (Z0) -- (1,1/3);
        \draw[black!60] (Z0) -- (1,1/2);
        \draw[black!60] (Z0) -- (1,2/3);
        \draw[black!60] (Z0) -- (1,3/4);
        \draw[black!60] (Z0) -- (1,4/5);
        \draw[black!60] (Z0) -- (1,10/11);
        \draw[black!60] (Z0) -- (1,1);

        % T_{1}: from (1,0) to {2}x(Q cap [0,1]) (sampled)
        \draw[black!60] (Z1) -- (2,0);
        \draw[black!60] (Z1) -- (2,1/10);
        \draw[black!60] (Z1) -- (2,1/4);
        \draw[black!60] (Z1) -- (2,1/3);
        \draw[black!60] (Z1) -- (2,1/2);
        \draw[black!60] (Z1) -- (2,2/3);
        \draw[black!60] (Z1) -- (2,3/4);
        \draw[black!60] (Z1) -- (2,4/5);
        \draw[black!60] (Z1) -- (2,10/11);
        \draw[black!60] (Z1) -- (2,1);

        \node[black!60!black] at (-0.55,0.55) {$T_{-1}$};
        \node[black!60!black] at (0.45,0.55)  {$T_{0}$};
        \node[black!60!black] at (1.45,0.55)  {$T_{1}$};
        \end{tikzpicture}}
        \end{center}

    Note that $T$ is path connected by a similar argument made in (a). However, it is not locally path connected. Indeed, any open ball around $(z,0)$ not only intersects the rays associated to $T_z$, but also intersects the rays associated to $T_{z-1}$. Hence there does not exist a path within any open ball from the point $(z,0)$ to $(z,q)$, where $q\in \bfQ \cap [0,1]$ is nonzero.
    \end{proof}


\end{document}